In conclusion, this project successfully designed, implemented, and utilized a new \inlinecode{clemgame}, Deal or No Deal, to conduct a rigorous evaluation of the negotiation capabilities of ten prominent LLMs. The findings reveal that while the latest generation of models, such as GPT-4.1, show a commendable ability to follow the game rules accurately and engage in superficially coherent negotiations, they still lack deep strategic reasoning. Models often fail to adapt their strategy to the specific goals of the game mode, falling back on generic negotiation patterns. Performance is highly variable across models, with some models clearly outperforming the rest.

Still, this work still has several limitations that could be potential avenues for future research. For example, only ten different models were tested in this project. A wider array of models could provide a more complete picture of the landscape, and reveal some underlying patterns explaining the differences in performance. Introducing human players to negotiate against LLMs could also provide a more realistic assessment of their capabilities. Further, all the game instances evaluated here were relatively small, i.e., few items and short dialogues, to ensure tractability. Performance might differ on more complex negotiation tasks. Another limitation of this work is the limited selection of languages. A broader linguistic analysis would therefore be beneficial. Future work could also explore the evaluation of the strictly competitive game mode. While its exclusion was a considered design choice based on methodological difficulties, it means the performance of the models in purely adversarial scenarios remains untested. The qualitative analysis performed for this project was also quite limited. Developing methods to automatically categorize the types of errors LLMs make, e.g., role confusion, logical contradiction, etc., would also enable more scalable qualitative analysis. Furthermore, future work could focus on a more fine-grained analysis of the negotiation strategies employed by the LLMs, perhaps by trying to identify specific tactics. Finally, this project used a single, fixed set of prompts. The impact of different prompting strategies on LLM performance was not explored. 

